\chapter{Overloads and coverage}
\label{ch:coverage}

Two features that belong together, because the second is what makes the first
readable: rosetta binds \textbf{overload sets}, and every generation writes a
machine-readable \code{coverage.json} saying what bound, what did not, and why.

\section{Overloads}

The reflection walk deduplicates member functions by \textbf{(name,
signature)}, so of

\begin{lstlisting}[style=cpp]
struct Grid {
    double at(int i) const;
    double at(int i, int j) const;
};
\end{lstlisting}

the whole set arrives in the IR. What happens next depends on the target
language, because ``two functions with one name'' is not something every runtime
can express.

\begin{center}
\small
\begin{tabular}{@{}llL{7.0cm}@{}}
\toprule
\textbf{Target} & \textbf{Policy} & \textbf{Why} \\
\midrule
\lang{python} & \textbf{all} & repeated \code{.def("at", \ldots)} builds one
  overload set; pybind dispatches on argument types \\
\lang{nanobind} & \textbf{all} & same model as pybind11 \\
\lang{julia} & \textbf{all} & each becomes a Julia method; multiple dispatch
  picks \\
\lang{qt} & \textbf{all} & builds an independent widget row per method ---
  nothing is keyed by name \\
\addlinespace
\lang{wasm} & first only & a second \code{.function("at", \ldots)} throws
  \code{BindingError} at module init \\
\lang{node} & first only & property descriptors are keyed by name; JS has no
  type dispatch \\
\lang{lua} & first only & \code{c["at"] = \ldots} is an assignment --- a second
  one overwrites \\
\lang{csharp} / \lang{java} & first only & the op table is a name-keyed map, and
  marshalling goes through JSON \\
\lang{typescript} & first only & the \code{.d.ts} describes the N-API module, so
  it must not promise what \lang{node} did not bind \\
\lang{rest} & first only & a route is a URL path; two overloads would register
  the same path \\
\lang{qml} & first only & \code{registerInvoker} is keyed by name, and QML calls
  with an untyped \code{QVariantList} \\
\bottomrule
\end{tabular}
\end{center}

``First'' always means \textbf{first-declared}, which is stable across
regenerations.

Every dropped overload appears in \code{coverage.json} with the reason
\code{overload\_not\_expressible}, so you can see exactly what a name-keyed
target left behind --- and reach it with an
\hyperref[sec:extensions]{extension method} under a distinct name if you need
it.

\begin{note}
This is also the one thing the runtime object model buys back. Through the
\lang{dynamic} projection, those same five name-keyed targets get their
overloads \emph{back}, because resolution happens at call time rather than at
generation time. See Chapter~\ref{ch:scriptable}.
\end{note}

\subsection{C++ name hiding is honoured}

\begin{lstlisting}[style=cpp]
struct Base    { int f(int) const; int f(int, int) const; };
struct Derived : Base { int f(double) const; };
\end{lstlisting}

\code{Derived} binds \textbf{only} \code{f(double)}. That is what a C++ caller
sees without a \code{using Base::f}, so binding the hidden base overloads would
hand scripts a call the C++ API does not offer. Both hidden declarations are
recorded as \code{hidden\_by\_derived} drops rather than vanishing.

A diamond-shared base member is still emitted exactly once. Two \emph{different}
bases declaring the same name both bind --- that call is ambiguous in C++
without qualification, but the binding has no ambiguity to resolve (each entry is
spliced from its own reflection), so binding both beats binding neither.

\section{\code{coverage.json}}

\code{generate()} writes it to \code{<out\_dir>/coverage.json} on every run,
always. An empty skip list is itself the useful answer, and a file that is
always present is one you can diff.

\begin{lstlisting}[style=json]
{
  "rosetta_coverage": 1,
  "reflection": [                       // never reached ANY backend
    { "class": "Shape",
      "dropped": [
        {"member": "operator<", "signature": "bool(const Shape &) const",
         "reason": "no_identifier"},
        {"member": "debug", "signature": "", "reason": "function_template"}
      ] }
  ],
  "targets": [
    { "target": "wasm",
      "bound": 3, "skipped": 1,
      "classes": [
        { "class": "Shape",
          "bound":   [ {"member": "area", "signature": "double () const"} ],
          "skipped": [ {"member": "scale", "signature": "void (double, double)",
                        "reason": "overload_not_expressible",
                        "detail": "the target binds methods by name and ..."} ] }
      ] }
  ]
}
\end{lstlisting}

\subsection{The two stages}

\textbf{\code{reflection}} lists members the walk never handed to any backend.
These are backend-independent:

\begin{center}
\small
\begin{tabular}{@{}lL{8.8cm}@{}}
\toprule
\textbf{Reason} & \textbf{Meaning} \\
\midrule
\code{no\_identifier} & an operator or conversion function --- no name to bind
  to \\
\code{function\_template} & a member template: no fixed parameter pack to
  splice \\
\code{hidden\_by\_derived} & a base overload hidden by a derived declaration of
  the name \\
\bottomrule
\end{tabular}
\end{center}

\textbf{\code{targets}} is what each backend's own gates decided:

\begin{center}
\small
\begin{tabular}{@{}lL{8.8cm}@{}}
\toprule
\textbf{Reason} & \textbf{Meaning} \\
\midrule
\code{unmarshalable\_signature} & no conversion for a type in the signature;
  \code{detail} names it \\
\code{overload\_not\_expressible} & the target keys methods by name and this one
  lost the tie \\
\code{sequence\_not\_adaptable} & a registered sequence with no
  \code{std::vector} boundary adapter here \\
\code{static\_callback\_receiver} & embind cannot give a static member the
  receiver a callback adapter needs \\
\code{extension\_has\_no\_member\_pointer} & this backend dispatches through a
  member pointer \\
\bottomrule
\end{tabular}
\end{center}

Both the \textbf{bound} and the \textbf{skipped} side are recorded. Counting
only skips cannot tell ``everything bound'' apart from ``this class never
reached this backend'' --- and that is exactly the distinction that matters when
a whole target silently produces nothing.

Free functions get their own \code{"functions": \{"bound": [\ldots], "skipped":
[\ldots]\}} array beside \code{"classes"}. On a library whose algorithms are
free functions over one data structure, that is most of the API.

\section{Using it}

The point is the \emph{diff}. Check \code{coverage.json} into your repository
next to the manifest, and a method that stops binding becomes a reviewable line
in a pull request instead of an \code{AttributeError} a week later.

To gate CI on it:

\begin{lstlisting}[style=sh]
jq -e '[.targets[].skipped] | add == 0' bindings/coverage.json
\end{lstlisting}

More usefully: allow known skips and fail on new ones by diffing against a
committed baseline. A binding surface that only ever grows deliberately is worth
the twenty lines of CI it costs.

\begin{warn}
One asymmetry to know before comparing numbers across targets:
\lang{typescript} and \lang{dynamic} record bound \emph{fields}, while the other
language backends record only methods. The top-level \code{bound} scalar is
therefore not comparable across those two groups --- count methods by signature
instead.
\end{warn}

\section{Instrumenting a backend of your own}

Skips are recorded at the existing gate sites, one line each:

\begin{lstlisting}[style=cpp]
for (const auto &m : k.methods) {
    if (!coverage::emit_overload(coverage::overloads::first_only,
                                 "my-backend", k, m)) {
        continue;                                   // records the dropped overload
    }
    if (!my_method_ok(m, c)) {
        coverage::note_skip("my-backend", k, m, "unmarshalable_signature",
                            my_skip_reason(m, c));  // a slug plus a sentence
        continue;
    }
    segs.push_back(my_method(k, m));
    coverage::note_bound("my-backend", k, m);
}
\end{lstlisting}

Free functions use \code{note\_bound\_function} / \code{note\_skip\_function},
fields \code{note\_bound\_field} / \code{note\_skip\_field}.

\begin{note}
Keep the reason function next to the predicate it explains. A reason that drifts
out of step with the decision is worse than no reason at all.
\end{note}

\section{A note for backend authors}

\code{\&T::name} names an \textbf{overload set}, not a function, so any emitter
that spells a member pointer for an overloaded name must cast it to the exact
signature. Use the shared helper:

\begin{lstlisting}[style=cpp]
const std::string mp = px_member_pointer(k, m);   // &T::f, or a static_cast
\end{lstlisting}

The gate is \code{GenMethod::is\_overloaded} --- does the \emph{C++ class}
overload the name? --- and \textbf{not} \code{overload\_count}, which is how
many entries reached the IR. A set whose siblings were all gated out still needs
the cast.
